\documentclass[10pt,a4paper]{amsart}
\usepackage[margin=19mm]{geometry}
\usepackage{amsmath,amssymb,mathtools}
\usepackage[scr=rsfs,cal=euler]{mathalfa}
\usepackage{xfrac}
\usepackage[unicode,hidelinks,bookmarks=false]{hyperref}
\newcommand{\ie}{i.e.\ }
\newcommand{\cf}{cf.\ }
\newcommand{\resp}{respectively\ }
\newcommand{\Subsection}{\subsection}
\providecommand{\nobreakdash}{\nobreak\hbox{-}}
\setlength{\emergencystretch}{2em}
\multlinegap0pt
\usepackage{bm}
\usepackage{bbm}
\usepackage{dsfont}
\usepackage{xspace}
\usepackage{cancel}
\usepackage{mathtools}
\usepackage{amsthm}
\makeatletter
\@ifundefined{prop}{\newtheorem{prop}{Prop}}{}
\@ifundefined{lemma}{\newtheorem{lemma}[prop]{Lemma}}{}
\makeatother
\newcommand{\C}[1]{\mathcal{#1}}
\newcommand{\OP}[1]{\operatorname{#1}}
\title[Directional asymptotic cones of groups equipped with bi-inva]{Directional asymptotic cones of groups equipped with bi-invariant metrics}
\author{Jarek K\k{e}dra, Assaf Libman}
\date{Source: arXiv:2207.12704v3 (CC BY 4.0)}
\begin{document}
\maketitle

\noindent\emph{Source and licence.} Directional asymptotic cones of groups equipped with bi-invariant metrics. Version arXiv:2207.12704v3 (CC BY 4.0), distributed under the Creative Commons Attribution 4.0 International licence, as recorded in the arXiv licence metadata for this exact version and shown on the abstract page https://arxiv.org/abs/2207.12704v3. Full rights notice: licenses/arxiv-2207.12704-CC-BY-4.0.txt.\par\smallskip

\noindent\textit{Editorial scope.} Counted unit: the Lemma whose source label is \texttt{L:non triviality product Theta} in the article (source0.tex lines 5631--5644, source label \texttt{L:non triviality product Theta}), delivered together with the article's own complete original proof of it (source0.tex lines 5647--5733, one \texttt{proof} environment of 87 lines). No mathematics is written, repaired or re-derived: statement and proof are the article's own text line for line. Required context retained uncounted: L:equal packets equal words (lines 5481--5485); L:small intervals reduction (lines 5509--5516); L:collision of partitions (lines 5576--5580). Every definition, display and label of those lines is retained unchanged. Omitted from this package and not redistributed: 1--5480, 5486--5508, 5517--5575, 5581--5630, 5645--5646, 5734--6425. Nothing inside a retained excerpt is omitted or altered: every retained body is delivered exactly as the article's lines are written, so no omission receipt of any kind accompanies this package. Citations: the retained excerpts carry no citation command at all, so no bibliography entry of the article is transcribed and no reference is redistributed. Reference adaptation: none was needed and none was made; every reference inside the delivered unit resolves inside it, so no unresolved reference or citation is produced. Printed slips kept literal: no printed word, symbol, hypothesis or number of the counted statement or of its proof was altered. Layout and class: the article's own class is \texttt{amsart} with packages including amsmath, amssymb, amscd, amsthm, indentfirst, amsfonts, enumitem, fontenc, hyperref, xy, stmaryrd, xcolor, xfrac, graphicx; the wrapper is the lane's pinned \texttt{amsart} editorial wrapper at 10pt on A4, which restates the theorem-like environments the retained text needs and adds the article's own macro definitions that the retained text needs (\texttt{\textbackslash C}, \texttt{\textbackslash OP}), with extra wrapper packages (bm, bbm, dsfont, xspace, cancel, mathtools). The delivered numbering is the unit's own. Assets: no retained span contains a figure environment, a graphics-inclusion command, a PDF-page inclusion, a table or a TikZ picture; no image file is copied into this package, the package declares no asset dependency and no image file is redistributed with it. Licence: version arXiv:2207.12704v3 of this article is distributed under the Creative Commons Attribution 4.0 International licence, as recorded in the arXiv licence metadata for this exact version and shown on the abstract page https://arxiv.org/abs/2207.12704v3. A full rights notice is delivered at licenses/arxiv-2207.12704-CC-BY-4.0.txt. Characters: every retained excerpt is pure ASCII, so the delivered engine, XeLaTeX with its default Latin Modern text font, reports no missing character.
\begin{lemma}\label{L:equal packets equal words}
Consider $\theta_1, \theta_2 \in \Theta$ and suppose that $\ell \geq \OP{lcm}(|\theta_1|,|\theta_2|)$.
Suppose that there exists a reduced word $v$ such that $|v| \geq \ell$ and such that $v$ is a $\theta_1$-packet and $v^{-1}$ is a $\theta_2$-packet.
Then $\theta_1=\theta_2$.
\end{lemma}

\begin{lemma}\label{L:small intervals reduction}
Consider cyclically reduced $g \in G$.
Then the removal of $m$ symbols from $g^n$ results in a  word whose reduced form is obtained from $g^n$ by the removal of at most $m$ intervals of total length at most
\[
\tfrac{|g| \cdot (1+|g|)}{\tau(g)} m.
\]
In particular it is a product of at most $m+1$ $g$-packets of total length at least $n|g|-\tfrac{|g|\cdot(1+|g|)}{\tau(g)} m$.
\end{lemma}

\begin{lemma}\label{L:collision of partitions}
Fix some numbers $\ell >0$ and $N>0$.
Suppose that $\C P_1=\{I_1,\dots,I_{m_1}\}$ and $\C P_2=\{J_1,\dots,J_{m_2}\}$ are partitions of $\{1,\dots,N\}$ into non-empty intervals where $N \geq \ell(m_1+m_2)$.
Then there are $I_i \in \C P_1$ and $J_j \in \C P_2$ which intersect at an interval of length at least $\ell$.
\end{lemma}

\begin{lemma}\label{L:non triviality product Theta}
Consider $\theta_1,\dots,\theta_p \in \Theta$  where $p \geq 1$.
Assume that $\theta_i \neq \theta_{i+1}$ for all $1 \leq i <p$.
Set $\ell=\OP{lcm}(|\theta_1|,\dots,|\theta_p|)$, and for every $1 \leq i \leq p$ set
\[
\kappa_i = \frac{|\theta_i|}{9\ell+|\theta_i| \cdot \tfrac{1+|\theta_i|}{\tau(\theta_i)}}.
\]
Let $n_1,\dots,n_p$ be integers such that $|n_i| >\tfrac{6 \ell}{|\theta_i|}$ for all $1 \leq i \leq p$.
Then
\[
\| \theta_1^{n_1} \cdots \theta_p^{n_p} \|_G >\min_{1 \leq i \leq p} |n_i| \cdot \kappa_i.
\]
% In particular $\theta_1^{n_1} \cdots \theta_p^{n_p} \neq 1$ in $F(A)$.
\end{lemma}

\begin{proof}
Let $g$ be the (not necessarily reduced) word $\theta_1^{n_1} \cdots \theta_p^{n_p}$.
Let $u$ be a sequence in $g$ of length 
\[
m \leq \min_{1 \leq i \leq p} |n_i| \cdot \kappa_i.
\] 
We will show that $g - u$ is not the trivial element and by this complete the proof since the norm on $G$ coincides with the cancellation norm.

Partition the sequence $u$ to subsequences $u_i \subseteq \theta_i^{n_i}$ of length $m_i$ each ($1 \leq i \leq p$).
Thus, $m=\sum_i m_i$.
By Lemma \ref{L:small intervals reduction} the removal of $u_i$ from $\theta_i^{n_i}$ yields, after reduction, a reduced word $w_i$ which is a product of at most $m_i+1$ $\theta_i$-packets of total length
\begin{align*}
|w_i| &\geq 
|n_i| \cdot |\theta_i| -m_i \cdot |\theta_i| \cdot \tfrac{1+|\theta_i|}{\tau(\theta_i)} 
\\
& \geq
|n_i| \cdot |\theta_i| - m \cdot |\theta_i| \cdot \tfrac{1+|\theta_i|}{\tau(\theta_i)} && \text{(since $m_i \leq m$)}
\\
& \geq
m \tfrac{|\theta_i|}{\kappa_i} - m \cdot |\theta_i| \cdot \tfrac{1+|\theta_i|}{\tau(\theta_i)} && \text{(since $m \leq \kappa_i |n_i|$)} \\
& =
9\ell m && \text{(by definition of $\kappa_i$)}.
\end{align*}
Then $g - u$ is equivalent to $w_1 \cdots w_p$. 
For any $1 \leq i <p$ we will now examine the reduction of $w_i w_{i+1}$ at the juncture.
This reduction boils down to a word $v_i$ such that $w_i=w_i'v_i$ and $w_{i+1}=v_i^{-1} w'_{i+1}$ and $w_i'w'_{i+1}$ is reduced.
We claim that for every $1 \leq i <p$
\begin{equation}\label{E:bound on length of v}
|v_i| \leq \min\left\{ \tfrac{|w_i|}{3},\tfrac{|w_{i+1}|}{3} \right\}+\ell.
\end{equation}
Assume false, i.e., $|v_i|>\min\left\{ \tfrac{|w_i|}{3},\tfrac{|w_{i+1}|}{3} \right\}+\ell$.
Then $|v_i|> (3m+1)\ell$.

If $m=0$ then $m_i=0$ for all $i$ and then $u$ is empty so $w_i=\theta_i^{n_i}$ for all $i$.
Then $v_i$ is a $\theta_i$-packet and $v_{i}^{-1}$ is a $\theta_{i+1}$ packet of length $>(3m+1)\ell \geq \ell \geq \OP{lcm}(|\theta_i|,|\theta_{i+1}|)$.
Lemma \ref{L:equal packets equal words}  implies that $\theta_i=\theta_{i+1}$ which is a contradiction.
This proves \eqref{E:bound on length of v} if $m=0$.
So we assume that $m\geq 1$.

Given $1 \leq i <p$, recall that $w_i$ is a product of at most $m_i+1$ $\theta_i$-packets and similarly $w_{i+1}$ is a product of at most $m_{i+1}+1$ $\theta_{i+1}$-packets.
Therefore $v_i$ is a product of at most $m_i+1$ $\theta_i$ packets and $v_i^{-1}$ is a product of at most $m_{i+1}$ $\theta_{i+1}$-packets.
%Then $v$ is a suffix of $w_i$ and $v^{-1}$ is a prefix of $w_{i+1}$. % where $|v| > \tfrac{9m \ell}{3}+\ell = 3m\ell+\ell$.
If we set $N=|v_i|$, then the partition of $w_i$ into the $\theta_i$-packets gives a partition $\C P_1$ of $\{1,\dots,N\}$ (the letters in $v_i$) into at most $m_i+1$ intervals.
Similarly, the partition of $w_{i+1}$ into the $\theta_{i+1}$-packets gives a partition $\C P_2$ of $\{1,\dots,N\}$ %(the letters in $v$) 
into at most $m_{i+1}+1$ intervals.
%In addition, since $m_i \neq 0$ or $m_{i+1} \neq 0$,
Then
% Since $m=m_i+m_{i+1} \geq 1$,
\[
%\begin{multline*}
N = 
|v_i| > 
(3m+1)\ell \geq %=
%3(m+1)\ell \geq
(m+2)\ell =
%3(m_i+m_{i+1})\ell+\ell \geq
%\\
%(m_i+m_{i+1})\ell+2(m_i+m_{i+1}+1)\ell \geq
%(m_i+m_{i+1}+2)\ell =
((m_i+1)+(m_{i+1}+1))\ell.
%\end{multline*}
\]
Lemma \ref{L:collision of partitions} implies that in the process of reduction of $w_iw_{i+1}$ there is a collision of a $\theta_i$-packet in $w_i$ with a $\theta_{i+1}$-packet in $w_{i+1}$ at an interval $v_i'$ in $v_i$ of length at least $\ell \geq \OP{lcm}(|\theta_i|,|\theta_{i+1}|)$.
Lemma \ref{L:equal packets equal words} implies that $\theta_i=\theta_{i+1}$ which is a contradiction. % (notice that $p \geq 2$ is necessary to obtain this contradiction).
This proves \eqref{E:bound on length of v} in the case $m \geq 1$.

Let us set $v_0$ and $v_p$ to be the empty words, which we regard as a prefix of $w_1$ and a suffix of $w_p$.
Clearly \eqref{E:bound on length of v} applies to them as well, where we agree that $w_0$ and $w_{p+1}$ are empty.
Then given $1 \leq i \leq p$
\[
|w_i|-|v_{i-1}|-|v_i| \geq
|w_i| - 2 \tfrac{|w_i|}{3} - 2\ell =
\tfrac{|w_i|}{3} - 2\ell.
\]
If $m=0$ then in particular $m_i=0$, in which case $w_i=\theta_i^{n_i}$, so
% By the assumption on $n_i$
\[
\tfrac{|w_i|}{3} - 2\ell  = \tfrac{|n_i| \cdot |\theta_i|}{3} - 2\ell > \tfrac{6\ell}{3}-2\ell =0.
\]
If $m>0$ then 
\[
\tfrac{|w_i|}{3} - 2\ell \geq \tfrac{9\ell m}{3}-2\ell \geq \ell >0.
\]
This proves that $v_{i-1}^{-1}$ and $v_i$ are disjoint intervals in $w_i$ for all $1 \leq i \leq p$, so $w_i=v_{i-1}^{-1}\widetilde{w_i}v_i$ for non empty subword $\widetilde{w_i}$.
It follows that the reductions at the junctures of $w_1 \cdots w_p$ do not affect each other so the reduced form of $w_1\cdots w_p$ is $\widetilde{w_1}\cdots \widetilde{w_p}$ which is non-trivial since $\widetilde{w_i}$ are not empty and reduced, and there may be no reductions at the junctures.
This completes the proof.
\end{proof}

\end{document}
